\section{Introduction}\label{sec:intro}

A controller that must keep a system inside a safe set, but cannot see the state, faces two questions. The first is qualitative: does the information available suffice, in principle, to guarantee safety? The second is quantitative: when it does not, how much safety survives? Viability theory answers the first question with set-valued objects---viability kernels, regulation maps and, under incomplete observation, obstruction certificates that prove an information state nonviable \cite{Aubin1991,Veliov1993,Doyen2000,Abaee2026calc}. Stochastic viability \cite{DeLaraDoyen2008,DoyenDeLara2010} and probabilistic reachability \cite{Abate2008} answer the second question under full observation, by a dynamic programme for the maximal probability of remaining in the constraint set. The partially observed, quantitative case sits between the two: the information state is a probability distribution over states, a belief, and the natural object is the \emph{belief-state safety value}
\[
V^{\Pi}_{k}(b)\;=\;\sup_{\pi\in\Pi}\;\Prob_{\pi}\bigl(x_t\in\calV\ \text{for } t=0,\dots,k\;\big|\;b\bigr),
\]
the maximal probability, within a declared policy class $\Pi$, of surviving $k$ steps from the belief $b$. Its value-one level is a viability statement; everything below it is measured by the deficit $1-V^{\Pi}_k(b)$.

This paper develops $V^{\Pi}_k$ as an exact theory. The word \emph{exact} is meant literally: every value, witness, threshold and count reported here is a rational number computed from the definitions in integer and rational arithmetic, and the structural theorems explain why the computations produce the numbers they do. The development has a deterministic layer, in which the kernels and observation maps are deterministic and the value measures how much prior mass a policy can keep alive, and a stochastic layer, in which sensing is noisy and the value is a genuine probability.

\paragraph{Contributions.}
\begin{enumerate}[leftmargin=1.6em,itemsep=1pt]
\item \emph{Exact sufficiency} (Theorem~\ref{thm:support}). On finite models with deterministic kernels and deterministic observation maps the support of the Bayesian posterior equals the set-valued post-state of the obstruction calculus; hence $V_k(b)=1$ if and only if $\supp b$ belongs to the feedback viable-set recursion $\Wfb_k$, and otherwise $1-V_k(b)\ge\min_{x\in\supp b}b(x)$. The criterion is necessary as well as sufficient, where the nearest output-feedback criterion of the viability literature \cite{Veliov1993} is sufficient only, and it admits history-dependent policies, where the necessary-and-sufficient criterion of \cite{Doyen2000} is confined to memoryless Lipschitz loops. The feedback and blind recursions differ, and the \emph{value of observation} $V^{\seq}_k-V^{\ol}_k$ is at most $1-\max_x b(x)$ whenever every state is individually survivable, a bound attained for every support size (Proposition~\ref{prop:feedback}).
\item \emph{The survivable-set antichain} (Theorem~\ref{thm:antichain}). Under deterministic kernels the $\alpha$-vectors of the value are exactly the indicators of the maximal jointly survivable subsets of the support. These form an antichain, so their number is at most $\binom{n}{\lfloor n/2\rfloor}$ for a support of size $n$ (Sperner), the bound is attained, and the deficit is $\min_S b(S^{c})$ over the maximal sets. This explains the three-type witness census of the certified instances and gives the value a finite, enumerable structure; the value stabilizes in the horizon after at most $2^{n}$ strict decreases (Proposition~\ref{prop:freeze}).
\item \emph{Certificates as deficit bounds} (Proposition~\ref{prop:deficit}). The chance-constrained common-action certificate, the minimal-mass certificate and horizon monotonicity are lower bounds on the exact deficit, attained on natural instances at symmetric and asymmetric priors.
\item \emph{Closed forms and the class lattice on the delayed hidden-regime class} (Theorems~\ref{thm:closed} and~\ref{thm:lattice}). For a hidden regime revealed after a blind window of length $\Tobs$, the hold-class, open-loop and fully observed safety values are obtained in closed form for all initial states, horizons, window lengths and prior weights. The inclusion of classes orders the values, both gaps are realized exactly and on disjoint cells---the value of flexibility on $\{2\le z_0<1+\min(k,\Tobs)\}$ and the value of observation on $\{1\le z_0<2\}$, each equal to $\min(w,1-w)$---and the classical timing obstruction is identified as a statement about the declared hold class (Corollary~\ref{cor:timing}). The complete certified value table, its $36/156$ census and its $384$ witnesses ($348$ distinct) are corollaries (Corollary~\ref{cor:grid}).
\item \emph{The price of late information} (Section~\ref{sec:deadline}). When a hidden drift parameter can only be learned by probing, a probe pinned to the deadline $\Tl$ shifts the viability threshold linearly, $z_0\ge\tfrac{21}{10}+\tfrac{\Tl}{10}$, between a flat free-probe benchmark and a probe-free one (Proposition~\ref{prop:deadline}); a general hold--probe--learn schedule has an exact threshold which reduces to an additive worst-loss law when the worst hold loss and the worst probe excursion are attained by a common parameter value, and not otherwise (Theorem~\ref{thm:additive}).
\item \emph{The stochastic layer} (Section~\ref{sec:stochastic}). Under noisy sensing the one-step deficit is the prior mass weighted by fatal-reading probabilities; a two-state sensor instance attains deficit $\min\{\varepsilon,b(x_1),b(x_2)\}$, so the minimal-mass bound is a deterministic-limit statement and not a law of the noisy layer. A single noisy probe of the hidden drift yields exactly the three value levels $\{0,1-\varepsilon,1\}$; majority reading over $n$ probes drives the wrong-majority probability below $(4\varepsilon(1-\varepsilon))^{n/2}$ at the cost of one unit of floor margin per additional reading---an exact exploration--safety exchange.
\item \emph{Robust interpolation} (Section~\ref{sec:dr}). A contamination ambiguity set of radius $\rho$ interpolates exactly between the expected and the worst-case reading at one stage, the adversary spending its whole budget on the worst outcome; the interpolation does not extend affinely across stages (a two-stage instance gives $\tfrac1{16}$ against $\tfrac18$), and a prior-contamination model has a closed form on the delayed class.
\end{enumerate}

\paragraph{Relation to the companion papers.}
The set-valued theory this paper quantifies is the obstruction calculus of \cite{Abaee2026calc}; its finite certification and worked-systems audits are treated in \cite{Abaee2026cert,Abaee2026worked}, and its minimax dual certificates in \cite{Abaee2026minimax}. Nothing proved there is re-derived differently here; the interface is Proposition~\ref{prop:deficit}, which reads each certificate as a lower bound on the exact deficit, and Corollary~\ref{cor:timing}, which scopes the timing certificate to its policy class. The two running instances---the delayed hidden-regime instance and the two-floor instance---are shared with the calculus so that the quantitative values can be compared cell by cell with the set-valued verdicts.

\paragraph{Organisation.}
Section~\ref{sec:setting} fixes the model, the policy classes, the safety value and its exact $\alpha$-vector representation. Section~\ref{sec:support} proves the support identity, the antichain theorem and stabilization. Section~\ref{sec:deficit} reads the certificates as deficit bounds. Section~\ref{sec:delayed} treats the delayed hidden-regime class: closed forms, the class lattice, the scoping of the timing obstruction and the certified grid. Section~\ref{sec:deadline} prices late information, Section~\ref{sec:stochastic} opens the stochastic layer, and Section~\ref{sec:dr} develops the contamination interpolation. Section~\ref{sec:related} positions the results, Section~\ref{sec:verification} describes the exact verification and the machine-checked layer, and Section~\ref{sec:discussion} discusses limitations and consequences. The supplementary material (S1--S8) carries the complete certified table, the asymmetric audit, the derivations of the jump loci, the learning-deadline tables, the exact probe-count values, the proofs for the contamination operator, the verification record and a cross-paper symbol table.
